\documentclass{article}
\usepackage[utf8]{inputenc}
\usepackage[margin=2cm,includefoot]{geometry}
\usepackage{booktabs}
\usepackage{graphicx}
\usepackage{amsmath, amssymb, amsthm}
\usepackage{physics}
\usepackage[shortlabels]{enumitem}
\usepackage{cancel}
\usepackage{lastpage}

\usepackage{hyperref}
\hypersetup{
    colorlinks,
    citecolor=black,
    filecolor=black,
    linkcolor=black,
    urlcolor=black
}

\usepackage{fancyhdr}
\pagestyle{fancy}
\lhead{L-head}
\rhead{R-head}
\chead{C-head}
\fancyfoot{}
\fancyfoot[R]{Page \thepage /\pageref{LastPage}}

\usepackage{listings}
\usepackage{color}
 \usepackage{titling}
 \usepackage{amsfonts}
 \usepackage{pdfpages}
 \usepackage{tikz}
 \usepackage{mathrsfs}
 \usepackage{bm}

\usepackage{calligra}

\DeclareMathAlphabet{\mathcalligra}{T1}{calligra}{m}{n}
\DeclareFontShape{T1}{calligra}{m}{n}{<->s*[2.2]callig15}{}
\newcommand{\scriptr}{\mathcalligra{r}\,}
\newcommand{\vbboldscriptr}{\pmb{\mathcalligra{r}}\,}

\newcommand{\vba}{\mathbf{A}}
\newcommand{\vbb}{\mathbf{B}}
\newcommand{\C}{\mathbf{C}}
\newcommand{\x}{\mathbf{\hat{x}}}
\newcommand{\y}{\mathbf{\hat{y}}}
\newcommand{\z}{\mathbf{\hat{z}}}
\newcommand{\vr}{\mathbf{\hat r}}
\newcommand{\thet}{\boldsymbol{{\hat \theta}}}
\newcommand{\ph}{\boldsymbol{{\hat \phi}}}
\newcommand{\ba}{\mathbf{a}}
\newcommand{\bb}{\mathbf{b}}
\newcommand{\vl}{\mathbf{l}}
\newcommand{\vv}{\mathbf{v}}

\newcommand{\R}{\mathbb{R}}
\newcommand{\p}{\mathbb{P}}
\newcommand{\N}{\mathbb{N}}
\newcommand{\Z}{\mathbb{Z}}
\newcommand{\Q}{\mathbb{Q}}
\newcommand{\I}{\mathbb{I}}
\newcommand{\E}{\mathbb{E}}

\newcommand{\mbeq}{\overset{!}{=}}
\newcommand{\ko}{\frac{1}{4\pi \epsilon_0}}

\newcommand{\PATH}{figures/}

\begin{document}
    \begin{center}
        \vspace*{1cm}
            
        \Huge
        \textbf{Solving a generalisation of tic-tac-toe or $k$-in-a-row game on infinite $\{m, n\}$ tilings of 2D space.}
            
        \vspace{0.5cm}
        \LARGE
        Subtitle
            
        \vspace{1.5cm}

        \large
        \textbf{Names}
            
         \vspace{0.5cm}
        description\\
        \vspace{0.5cm}
        University

    \end{center}
\newpage
\tableofcontents
\newpage

\section{Introduction}
\subsection{Insights}
In $\{m, n\}$ tiling let 
\begin{align}
    m = 2i: i\in \N. \label{m_condition}
\end{align}
It matters so that there ale straight lines going through the shapes. Then in $n = 2j + 1$ tilings, diagonals do not make sense. In $n = 2j$ tilings, however, diagonals make sense. We then distinguish between a $\{m, n = 2j\}_m$ and $\{m, n = 2j\}_{vn}$ configurations named after the Moore (M) and Von Neumann (VN) neighbourhoods. The Moore version permitting diagonals to be used and the Von Neumann version only permitting edge-connected series. The classical tic-tac-toe is then played on a $\{4, 4\}_m$ board. In the spirit of $m,n,k$ games we include also the the length of the winning sequence as follows: $\{m, n\}^k$. The number of axis that you can win on then becomes $m$ in the M configuration and $m/2$ in the VN configuration.

\subsection{Notation}

\subsubsection{Schläfli symbol}

The Schläfli symbol is used to compactly denote any regular polytope or tessellation. For our case it will be of the form $\{p, q\}$, where $\{p\}$ is a regular $p$-sided polygon and $q$ is the number of these polygons joined at each vertex.

\subsubsection{Game versions}
$(1, 1, 1)$ where each player occupies 1 position in a move. $(1, 2, 2)$ where the first move occupies 1 position and then both players occupy 2 positions per move.

\subsection{Maker-Breaker games}
In a Maker-Breaker game, the Maker is the player trying to achieve the winning objective and Breaker is the player trying to stop the Maker from achieving this.

\subsubsection{Insights into Maker-Maker games}
If a $k$-game is a Breaker win, we know that the original Maker-Maker version is a draw.

\subsubsection{Strategy stealing argument for $(1, 1, 1)$}
In a Maker-Maker game, if there was a strategy such that P2 would have a winning strategy, P1 could just steal the strategy in a way where P1 would make a random move and then effectively play the role of P2 and employ that winning strategy. This is a contradiction as both P1 and P2 have a winning strategy. Therefore only P1 can have a winning strategy.

\subsubsection{Pairing strategies}
A pairing strategy is a strategy for Breaker where Breaker pairs up positions on the board in such a way that if Maker occupies one position in a pair, Breaker can immediately occupy the other position in the pair and thus prevent Maker from winning. If Breaker can successfully create such a pairing strategy that covers all possible winning combinations for Maker, then Breaker has a winning strategy.

\section{Tiling spaces}
No spherical tilings since they are all finite and do not form an infinite playing board.

The proof for any tiling with $k = 1$ is trivial as the first player wins by making any move.

\subsection{Euclidean tilings}
The only regular tilings of Euclidean space are $\{2, \infty\}, \{3, 6\}, \{4, 4\}, \{6, 3\}, \{\infty, 2\}$. We do not look at the $\{\infty, 2\}$ tiling as it does not provide an infinite playing board. We also exclude the $ \{3, 6\}$ tiling as it does not satisfy condition \ref{m_condition}. Finally, we also do not spend too much time looking at the $\{4, 4\}$ tiling as this is the case with regular $m, n, k$-games and has been studied extensively. We are left with the $\{2, \infty\}$ and $\{6, 3\}$ tilings, which we take a close look at.

\subsubsection{The $\{2, \infty\}$ tiling}
This is a very trivial case as it is a win for P1 in the case of $k \le 2$ and a draw in the case $k > 2$.
\par
For the case $k = 2$ P1 has a winning strategy. We can see that the $\{2, \infty\}$ has a reflective symmetry. Therefore the is only one game (up to symmetry) of $\{2, \infty\}^2$. It is shown in Figure \ref{fig:{2, infty}-2}.

\begin{figure}[ht]
    \centering
    \includegraphics[width=0.5\linewidth]{\PATH{2, infty}^2_win.png}
    \caption{Drawing of the winning strategy for $\{2, \infty\}^2$.}
    \label{fig:{2, infty}-2}
\end{figure}

\par
For the case $k > 2$ We can show this using a simple pairing argument where we tile the space as in Figure \ref{fig:{2, infty}}.
\begin{figure}[ht]
    \centering
    \includegraphics[width=0.5\linewidth]{\PATH{2, infty}_pair.png}
    \caption{Drawing of the pairing strategy for $\{2, \infty\}^3$.}
    \label{fig:{2, infty}}
\end{figure}\\
\textbf{NEEDS PROOF OF SAME THING FOR $(1, 2, 2)$}

\subsubsection{The $\{6, 3\}$ tiling}
Since the vertex degree in this tiling is odd, we will not consider diagonals.

\begin{figure}[ht]
    \centering
    \includegraphics[width=0.5\linewidth]{\PATH 6_3_7-tiling.png}
    \caption{Drawing of the pairing strategy for $\{6, 3\}^7$.}
    \label{fig:{6, 3}-7}
\end{figure}

\begin{figure}[ht]
    \centering
    \includegraphics[width=0.5\linewidth]{\PATH 6_3-tile.png}
    \caption{One repeating tile in the drawing $\{6, 3\}^7$ tiling.}
    \label{fig:{6, 3}-7-tile}
\end{figure}

Figure \ref{fig:{6, 3}-7} shows a winning breaker strategy for the $\{6, 3\}^7$ game. The lines represent pairs like in a regular pairing strategy.
The rest is tiled by the 4x3 rhombuses. A zoomed in version of one of these tiles is in Figure \ref{fig:{6, 3}-7-tile}.
The couloured lines represent all the directions in which breaker has to play at least once to win.\\
This has been shown to be a draw by T. G. L. Zetters \cite{guy_selfridge_zetters_1980_s10}.\\
\textbf{NEEDS PROOF FOR $k < 5$.}\\
\begin{figure}[ht]
    \centering
    \includegraphics[width=0.5\linewidth]{\PATH{6, 3}_win4.png}
    \caption{Forced Winning configuration (up to symmetry) for $\{6, 3\}^4.$}
    \label{fig:6, 3 - win 4}
\end{figure}

\begin{figure}[ht]
    \centering
    \includegraphics[width=0.5\linewidth]{\PATH{6, 3}_last_moves-4.png}
    \caption{Threat illustration for $\{6, 3\}^4$.}
    \label{fig:6, 3 - threats 4}
\end{figure}

\textbf{NEEDS PROOFS FOR $(1, 2, 2)$}

\subsection{Hyperbolic tilings}
There is an infinite amount of hyperbolic tilings. We will therefore look at some representative examples and then draw general conclusions.

\newpage
\bibliographystyle{plain}
\nocite{*}
\bibliography{bibliography}
\end{document}